\begin{theorem}
\label{thm:5term}
Partition $\{1,\ldots n\}$ into three nonempty sets: $A$, $B$, $C$, where $|C|=r$. Then 
\begin{equation}
\label{eq:5term}
\big[(A\cup B \mid C)\big]_r=\sum_{S\subseteq C}(-1)^{|S|}\big[( A\cup S \mid B\cup (C\setminus S) ) \big]_r
\end{equation}
\end{theorem}

To illustrate Theorem~\ref{thm:5term}, let $r=3$ and suppose $C = \{c_1, c_2,c_3\}$. Figure~\ref{fig:recurrence} gives a diagrammatic representation of the $9$-term recurrence:
{\small
\begin{align*}
\big[&(A\cup B \mid C)\big]_3
=\\ 
&\big[(A \mid B \cup C)\big]_3
-\big[(A\cup \{c_1\}  \mid B \cup \{c_2,c_3\})\big]_3
-\big[(A\cup \{c_2\}  \mid B \cup \{c_1,c_3\})\big]_3\\
-&\big[(A\cup \{c_3\}  \mid B \cup \{c_1,c_2\})\big]_3 
+\big[(A\cup \{c_1,c_2\} \mid B  \cup \{c_3\})\big]_3
+\big[(A\cup \{c_1, c_3\}  \mid B \cup \{c_2\})\big]_3\\
+&\big[(A\cup \{c_2,c_3\} \mid B \cup \{c_1\})\big]_3
-\big[(A\cup C \mid B)\big]_3
\end{align*}
}

\begin{figure}[htbp]
\begin{center}
\includegraphics[width=\textwidth]{figures/9termrecurrencefinal.pdf}
\caption{An illustration of the recurrence of Theorem~\ref{thm:5term} in the case $r=3$.}
\label{fig:recurrence}
\end{center}
\end{figure}

\begin{proof}[Proof of Theorem~\ref{thm:5term}]
        We prove Theorem~\ref{thm:5term} using Pl\"ucker coordinates and relations. As a shorthand, for $W$ any set of positive integers, we write $W' = \{w + n : w \in W\}$.
        
    First, we rewrite the left side of \eqref{eq:5term} in terms of Pl\"ucker variables, using Equation \eqref{eq:PlutoMinor}.  By definition, there is only one jellyfish tableau for the ordered set partition $(A\cup B | C)$, since $|C|=r$. Call this tableau $\widehat{T}$.
    We have
    \begin{equation}\label{eq:expand_the_jellyfish_inv}
    \begin{split}
            \big[(A\cup & B \mid C)\big]_r = \sum_{T\in \rstab_r(A\cup B \mid C)}\sgn(T) \; \rspoly(T)  \\
          &=  \sgn(\widehat{T}) \;  M_{R_{A \cup B}(\widehat{T})}^{A \cup B} \cdot M_{R_C(\widehat{T})}^{C} \\
          &=  \sgn(\widehat{T}) \; (-1)^{|A \cup B|} \Delta_{([n]\setminus R_{A \cup B}(\widehat{T}))  A' B'}(\Phi(M)) \cdot (-1)^{|C|}\Delta_{([n]\setminus R_C(\widehat{T}))  C'}(\Phi(M))\\
          &= (-1)^n \sgn(\widehat{T}) \;  \Delta_{[n-r + 1, n]  A' B'}(\Phi(M)) \cdot \Delta_{[r+1, n]  C'}(\Phi(M))
    \end{split}
      \end{equation}
Here, the first and second equalities are by definition, while the third is using \eqref{eq:PlutoMinor} and the fourth line collects the sign.

We will apply the classical Pl\"ucker relation in a form that we borrow from \cite[Eqn.~9.1]{Fulton:YoungTableaux}:
\begin{equation}\label{eq:Fulton}
\Delta_{i_1,i_2,\ldots,i_d}\cdot\Delta_{j_1,j_2,\ldots,j_d}=\sum\Delta_{i_1',i_2',\ldots,i_d'}\cdot\Delta_{j_1',j_2',\ldots,j_d'},
\end{equation}
where the sum is over all pairs obtained by interchanging a fixed set of $k$ of the subscripts $j_1,j_2,\ldots,j_d$ with $k$ of the subscripts in $i_1,i_2,\ldots,i_d$, maintaining the order in each. Note that whenever the resulting subscripts are rearranged, signs are introduced. In our setting, we will have $k = |B|$. 

Before applying the Pl\"ucker relation, we rearrange the last line of \eqref{eq:expand_the_jellyfish_inv} in order to better match \eqref{eq:Fulton}, obtaining:
\begin{equation}\label{eq:expand_the_jellyfish_inv_some_more}
\begin{split}
        \big[(A\cup B \mid C)\big]_r 
        &= (-1)^n \sgn(\widehat{T}) \;   \Delta_{[r+1, n]  C'}(\Phi(M)) \cdot \Delta_{[n-r + 1, n]  A'B'}(\Phi(M)) \\
        &= (-1)^n \sgn(\widehat{T}) \;   \Delta_{[n-r+1, n]  C'  [r+1,n-r]}(\Phi(M)) \cdot \Delta_{[n-r + 1, n]  A' B'}(\Phi(M)). 
        \end{split}
        \end{equation}
In the second equality above, we also moved the interval $[r+1,n-r]$ to the end of its Pl\"ucker coordinate to simplify things later. This move does not affect the global sign because each element in $[r+1,n-r]$ moves past the $r$ elements of $[n-r+1,n]$ and the $r$ elements of $C'$, for an even total number of swaps.

The particular Pl\"ucker relation we apply is the one in which $B'$ is the fixed set of $k$ subscripts that gets interchanged. What happens is that the first $\ell$ indices in $B'$ interchange with a subset $Z$ of $[r+1,n]$ and the rest interchange with a subset $S'$ of $C'$. We can match each term from the right side of \eqref{eq:Fulton} with the column-permuted jellyfish tableau that is obtained from $\widehat{T}$ by moving the elements of $B$ from column $1$ to column $2$, moving the elements of $S$ from column $2$ to column $1$, and so that the elements of $Z$ are the row indices of the empty boxes in column $1$. By the definition of jellyfish tableaux, this also determines the row indices of the empty boxes in column $2$. See Figure~\ref{fig:jellyfish_sign_calc} for an example of this construction (the unexplained notations in that figure will be introduced later for a sign calculation). By Lemma~\ref{lem:row_col_swap}, the ordering of the labels within each column is immaterial; however, the ordering we impose here will be useful later. Thus, each choice of subsets $S$ and $Z$, corresponds to a unique jellyfish tableau contributing to the right hand side of \eqref{eq:5term} by sorting the columns in increasing order. Conversely, one can check that each of the jellyfish tableaux corresponds to a term of the Pl\"ucker relation.

\begin{figure}[htb]
    \centering
    \begin{tikzpicture}
    \node (T1) {\ytableaushort{{A_1}Y,{B_1}S,a,a,b,a,b,b,a,b}};
    \node[right= 18em of T1] (T2) {\ytableaushort{{A_1}Y,S{B_1},a,a,\none b,a,\none b, \none b,a, \none b}};
    \node[right=4.5em of T1] (fake1) {};
    \node[right=7em of fake1] (fake2) {};
    \node[below left = 3.2em and 3em of T2]  (fake3) {\phantom{123}};
    \node[below = 3em of T1] (P1) {$\Delta_{[n-r+1,n] Y'S'[r+1,n-r]} \Delta_{[n-r+1,n]A_1'B_1'(A_2' \shuffle B_2')} $};
    \node[below right = 0.5em and -4em of P1] (P2) {$\Delta_{[n-r+1,n]Y'B'_1(Z^C \shuffle B'_2)} \Delta_{[n-r+1,n]A'_1S'(A'_2 \shuffle Z)}$};
    \draw[thick, ->] (fake1) -- (fake2);
    \draw[thick, ->] (P1) -- (fake3);
    \end{tikzpicture}
    \caption{An illustration of the proof of Theorem~\ref{thm:5term}.}
    \label{fig:jellyfish_sign_calc}
\end{figure}

Theorem~\ref{thm:5term} is then proved, except for comparing the signs from the jellyfish tableaux with the signs arising from sorting indices in \eqref{eq:Fulton}. 

As a warm up, we first consider the case $|A \cup B| = r$. Note that in this case, $n=2r$, so $[r+1,n-r]=\emptyset$.  We can assume our  Pl\"ucker monomial from \eqref{eq:expand_the_jellyfish_inv_some_more} looks like $\Delta_{[n-r+1,n]  Y'S'}\Delta_{[n-r+1,n]  A'B'}$, where $Y = C \setminus S$. Since the elements of the subscripts on these Pl\"ucker variables are out of order, there is an associated sign: $-1$ to the number of inversions. We want to see how this sign changes when we rearrange indices to obtain $\Delta_{[n-r+1,n]  Y'B'} \Delta_{[n-r+1,n]  A'S'}$.  Note that for each choice of $S$, we are choosing a different ordering in the initial Pl\"ucker monomial from \eqref{eq:expand_the_jellyfish_inv_some_more}, but that this is acceptable, since we only need to understand the relation between the signs of the two monomials of interest. In comparing the signs of these monomials, it is straightforward to see that the signs differ by $-1$ to the number of pairs $(t_1, t_2) \in (B \cup S) \times (Y \cup A)$ with $t_1 < t_2$.

We now perform the corresponding sign calculation for the corresponding jellyfish tableaux.
By Lemma~\ref{lem:row_col_swap}, we can instead compare the column-permuted jellyfish tableaux 
\ytableaushort{A Y, B S} and \ytableaushort{A Y, S B} , where the elements of each set are arranged vertically in a column in some fixed order. Since inversions in a column do not contribute to the sign of a jellyfish tableau, here there is a sign change for each pair $(t_1, t_2) \in (B \cup S) \times (Y \cup A)$ with $t_1 < t_2$, matching the sign changes from shuffling indices in the Pl\"ucker monomials, as described in the previous paragraph; in addition, there is a sign change for each of the bottom $|B| = |S|$ rows, as desired to match the sign appearing on the right-hand side of \eqref{eq:5term}. This completes the proof in the special case $|A \cup B| = r$.

Now we consider the general case of this sign calculation. 
Let $B_1$ be the first $|S|$ elements of $B$ and let $B_2 = B \setminus B_1$. Similarly, let $A_1$ be the first $r - |S|$ elements of $A$ and let $A_2 = A \setminus A_1$.

We study the column-permuted version of $\widehat{T}$ for which the left column has $A_1$ (in increasing order) above $B_1$ (in increasing order) occupying the body of the jellyfish, and then $A_2$ and $B_2$ in the following rows, shuffled together so that $B_2$ occupies the rows indexed by $Z$. Write $A_2 \shuffle B_2$ for this ordering of $A_2 \cup B_2$. We also write $Z^C$ for $[r+1,n-r] \setminus Z$. The right column of our column-permuted $\widehat{T}$ has $Y = C \setminus S$ (in increasing order) above $S$ (in increasing order). See the left tableau of Figure~\ref{fig:jellyfish_sign_calc} for an illustration.

We can assume our corresponding Pl\"ucker monomial from \eqref{eq:expand_the_jellyfish_inv_some_more} is written in a similar way, as \[
\Delta_{[n-r+1,n]  Y'S' [r+1,n-r]} \Delta_{[n-r+1,n]  A_1'B_1' (A_2' \shuffle B_2')},
\]
where $Y = C \setminus S$. Then we want to compare the sign of this with the sign for
\[
\Delta_{[n-r+1,n]Y'B'_1(Z^C \shuffle B'_2))} \Delta_{[n-r+1,n]A'_1S'(A'_2 \shuffle Z))}.
\]
Again, note that for each choice of $S$, we are choosing a different ordering in the corresponding Pl\"ucker monomial. 
In comparing the signs of these monomials, it is again straightforward to see that swapping $B_1'$ with $S'$ results in a sign difference of $-1$ to the number of pairs $(t_1, t_2) \in (B_1 \cup S) \times (Y \cup A_1)$ with $t_1 < t_2$ together with $-1$ to the number of pairs $(t_1, t_2) \in (B_1 \cup S) \times A_2$ with $t_1 > t_2$. 
Furthermore, the sign difference in swapping $B_2'$ and $Z$ comes from the fact that all the elements of $B_2'$ are larger than the elements of $Z^C$ and all the elements of $Z$ are smaller than the elements of $A_2'$. So after the swap, all the elements of $B_2'$ have inversions with  the elements of $[r+1,n-r]$ to their right, while all the elements of $Z$ have inversions with the elements of $A_2'$ to their left. Since the positions occupied by $B_2'$ and $Z$ are corresponding, this results in a total of $n-2r$ inversions. Since $2r$ is even, this results in a contribution of $(-1)^n$.

We now compute the corresponding sign difference between the corresponding column-permuted jellyfish tableaux. See Figure~\ref{fig:jellyfish_sign_calc} for a visual comparison of these tableaux. It is again straightforward to see that swapping $B_1$ with $S$ results in a sign difference of $-1$ to the number of pairs $(t_1, t_2) \in (B_1 \cup S) \times (Y \cup A_1)$ with $t_1 < t_2$ together with $-1$ to the number of pairs $(t_1, t_2) \in (B_1 \cup S) \times A_2$ with $t_1 > t_2$. Furthermore, there is a sign change for each of the bottom $|B_1| = |S|$ rows.

So, continuing the calculation from \eqref{eq:expand_the_jellyfish_inv_some_more}, we have that 
\begin{align*}
   [(A\cup B \ | \ C)]_r =& (-1)^n\sgn(\widehat{T})\sum_{S\subseteq C}\sgn(\widehat{T})(-1)^n(-1)^{|S|}\big[( A\cup S \mid B\cup (C\setminus S) ) \big]_r \\
    =& \sum_{S\subseteq C}(-1)^{|S|}\big[( A\cup S \mid B\cup (C\setminus S) ) \big]_r
\end{align*}
as desired.
\end{proof}
